\documentclass[12pt]{article}
\usepackage[utf8]{inputenc}
\usepackage[T1]{fontenc}
\usepackage{hyperref}
\title{Dayao Yi Language Resources for K-12 Curriculum Development}
\author{Compiled by Claude Sonnet 5.5 (claude-sonnet-5-5) for Andr\'as Kornai}
\date{}
\begin{document}
\maketitle

\section{Introduction}
Dayao Yi (ISO 639-3 yio; Glottocode daya1245) is a Yi (Ngwi/Loloish) variety of Dayao County, Yunnan, China. Glottolog lists it under the ``Bookkeeping'' label rather than in a real family, which reflects classification uncertainty. Glottolog reports David Bradley's fieldwork finding two Lipo dialects to be mutually intelligible. The sources do not settle whether Dayao Yi is the same as Lolopo (Central Yi), spoken in Dayao and Nanhua Counties. I found no speaker count and no description specific to Dayao Yi.

\section{Primary Educational Resources}
None specific to Dayao Yi were found.

\section{Linguistic and Reference Materials}
\begin{itemize}
\item Glottolog entry for Dayao Yi (yio). \url{https://glottolog.org/resource/languoid/iso/yio}
\item Wikipedia, ``Yi people'': places Dayao Lipo in a Central Yi group. \url{https://en.wikipedia.org/wiki/Yi_people}
\item Wikipedia, ``Dayao County'': background on the region, including local dialect names. \url{https://en.wikipedia.org/wiki/Dayao_County}
\end{itemize}

\section{Religious and Community Resources}
None found.

\section{Digital and Audio Resources}
\begin{itemize}
\item Omniglot, Lolopo (Central Yi): a related variety, written in the Yi script; it may overlap with Dayao Yi and is not confirmed to be the same language. \url{https://www.omniglot.com/writing/lolopo.htm}
\end{itemize}

\section{Curriculum Development Context}
Nothing specific to Dayao Yi was found. Any work would likely draw on resources for Central Yi / Lolopo and the standard Yi script taught in China, after confirming with local speakers and Yunnan education authorities how closely those match the Dayao variety.
\end{document}